\documentclass[11pt,a4paper]{article}

% --- MOTEUR ET POLICES ---
\usepackage{fontspec}

% --- MATHÉMATIQUES ---
\usepackage{amssymb}
\usepackage{amsthm}

% --- BOÎTES ET MISE EN PAGE ---
\usepackage[skins,breakable]{tcolorbox}
\usepackage{geometry}
\geometry{margin=2.5cm, headheight=20pt}

% --- LANGUE ---
\usepackage{polyglossia}
\setdefaultlanguage{french}

% --- THÉORÈMES ET ENVIRONNEMENTS ---
% Compteur attaché aux sections
\newtheorem{theorem}{Théorème}[section]

% Style tcolorbox pour l'environnement theorem
\tcolorboxenvironment{theorem}{
    breakable,
    enhanced,
    colback=white,
    colframe=black!75,
    fonttitle=\bfseries,
    boxrule=0.8pt,
    arc=2pt,
    before skip=10pt,
    after skip=15pt
}

\begin{document}

\section{Analyse réelle}

\begin{theorem}[Lemme de la corde universelle]
Soit $f : [0,1] \to \mathbb{R}$ une fonction continue telle que $f(0) = f(1)$. 
Alors il existe $c \in [0, 1/2]$ tel que $f(c) = f\left(c + \frac{1}{2}\right)$.
\end{theorem}

\begin{proof}
\noindent
\textbf{1. Définition d'une fonction auxiliaire.} \\
Posons $g : [0, 1/2] \to \mathbb{R}$ définie par :
\[ g(x) = f(x) - f\left(x + \frac{1}{2}\right) \]

\textbf{2. Continuité de $g$.} \\
La fonction $f$ est continue sur $[0,1]$ par hypothèse. La fonction $x \mapsto x + \frac{1}{2}$ est continue sur $[0, 1/2]$ et à valeurs dans $[1/2, 1] \subset [0,1]$. Par composition, $x \mapsto f\left(x + \frac{1}{2}\right)$ est continue sur $[0, 1/2]$. Par somme, $g$ est continue sur $[0, 1/2]$.

\textbf{3. Évaluation aux bornes.} \\
Calculons $g(0)$ et $g(1/2)$ :
\[ g(0) = f(0) - f\left(\frac{1}{2}\right) \]
\[ g\left(\frac{1}{2}\right) = f\left(\frac{1}{2}\right) - f(1) \]
En sommant ces deux quantités, on obtient :
\[ g(0) + g\left(\frac{1}{2}\right) = f(0) - f(1) \]

\textbf{4. Utilisation de l'hypothèse $f(0) = f(1)$.} \\
D'après l'hypothèse, $f(0) - f(1) = 0$, donc :
\[ g(0) + g\left(\frac{1}{2}\right) = 0 \implies g\left(\frac{1}{2}\right) = -g(0) \]

\textbf{5. Application du Théorème des Valeurs Intermédiaires (TVI).} \\
Deux cas se présentent :
\begin{itemize}
    \item \textbf{Cas 1 :} Si $g(0) = 0$, alors $f(0) = f(1/2)$. Le réel $c = 0 \in [0, 1/2]$ convient.
    \item \textbf{Cas 2 :} Si $g(0) \neq 0$, alors $g(0)$ et $g(1/2)$ sont non nuls et de signes opposés. Puisque $g$ est continue sur le segment $[0, 1/2]$, d'après le Théorème des Valeurs Intermédiaires, il existe $c \in \, ]0, 1/2[ \subset [0, 1/2]$ tel que $g(c) = 0$.
\end{itemize}

\textbf{6. Conclusion.} \\
Dans les deux cas, il existe $c \in [0, 1/2]$ tel que $g(c) = 0$, c'est-à-dire $f(c) = f\left(c + \frac{1}{2}\right)$.
\end{proof}

\end{document}